%%%%%%%%%%%%%%%%%%%%%%%%%%%%%%%%%%%%%%%%%%%%%%%%%%%%%%%%%%%%%%%%%%%%%%%%%%%%%%%%
%   Copyright 2009, Oracle and/or its affiliates.
%   All rights reserved.
%
%
%   Use is subject to license terms.
%
%   This distribution may include materials developed by third parties.
%
%%%%%%%%%%%%%%%%%%%%%%%%%%%%%%%%%%%%%%%%%%%%%%%%%%%%%%%%%%%%%%%%%%%%%%%%%%%%%%%%

\chapter{Convenience Functions and Types}
\chaplabel{lib:convenience}

\section{Convenience Functions}
\seclabel{convenience-functions}

%tuple[\T extends Tuple\](x: T): T

\Method{\(\VAR{tuple}\llbracket{}T \KWD{extends}\mskip 4mu plus 4mu\TYP{Tuple}\rrbracket(x\COLON T)\COLON T\)}

The function \VAR{tuple} returns its argument as a tuple expression.

%coerce_[\T\](x: T): T

\Method{\(\mathrm{coerce}\llbracket{}T\rrbracket(x\COLON T)\COLON T\)}

The function \EXP{\mathrm{coerce}} returns its argument as the given type.


\section{Convenience Types}
\seclabel{convenience-types}

An optional value \VAR{v} of type \EXP{\TYP{Maybe}\llbracket{}T\rrbracket}
is either \EXP{\TYP{Nothing}\llbracket{}T\rrbracket} or
\EXP{\TYP{Just}(v)} declared as follows:
%(* Optional Values *)
%trait Maybe[\T\] comprises { Nothing[\T\], Just[\T\] }
%  isNothing: Boolean
%end
%value object Nothing[\T\] extends Maybe[\T\]
%end
%object Just[\T\](just: T) extends Maybe[\T\]
%end
\begin{Fortress}
\(\mathtt{(*}\;\hbox{\rm  Optional Values \unskip}\;\mathtt{*)}\)\\
\(\KWD{trait}\:\TYP{Maybe}\llbracket{}T\rrbracket \KWD{comprises} \{\,\TYP{Nothing}\llbracket{}T\rrbracket, \TYP{Just}\llbracket{}T\rrbracket\,\}\)\\
{\tt~~}\pushtabs\=\+\(  \VAR{isNothing}\COLON \TYP{Boolean}\)\-\\\poptabs
\(\KWD{end}\)\\
\(\KWD{value}\;\;\KWD{object}\:\TYP{Nothing}\llbracket{}T\rrbracket \KWD{extends}\:\TYP{Maybe}\llbracket{}T\rrbracket\)\\
\(\KWD{end}\)\\
\(\KWD{object}\:\TYP{Just}\llbracket{}T\rrbracket(\VAR{just}\COLON T) \KWD{extends}\:\TYP{Maybe}\llbracket{}T\rrbracket\)\\
\(\KWD{end}\)
\end{Fortress}
The empty case is one object for each type \VAR{T},
and each use writes the type,
as in \EXP{\TYP{Nothing}\llbracket\TYP{String}\rrbracket}.
The following declaration of a single \TYP{Nothing} for every
instantiation of \TYP{Maybe} is not allowed:
%(* Not allowed! *)
%object Nothing extends Maybe[\T\] excludes Just[\T\] where {T extends Object}
%end
\begin{Fortress}
\(\mathtt{(*}\;\hbox{\rm  Not allowed! \unskip}\;\mathtt{*)}\)\\
\(\KWD{object} \TYP{Nothing} \KWD{extends} \TYP{Maybe}\llbracket{}T\rrbracket \KWD{excludes} \TYP{Just}\llbracket{}T\rrbracket \KWD{where} \{T \KWD{extends} \TYP{Object}\}\)\\
\(\KWD{end}\)
\end{Fortress}
It would make \TYP{Nothing} a \EXP{\TYP{Maybe}\llbracket{}T\rrbracket}
for every \VAR{T} at once,
which instantiation exclusion forbids (\secref{where-clauses});
and an object declaration has no \KWD{excludes} clause
(\secref{object-decls}), an object trait type already excluding
every type that is not its supertype (\secref{object-traits}).
\revision{revival-nothing}{The Working Draft of February 2011 declared
\TYP{Nothing} with the declaration marked ``Not allowed''.
The declaration of \EXP{\TYP{Nothing}\llbracket{}T\rrbracket} is the
library's, whose own source says it
``will become a non-parametric singleton when we get where clauses working'';
instantiation exclusion closes that road.}
